\chapter{Methodology}
\label{ch:methodology}

This chapter specifies the experimental  of the thesis. We first
formalize the MEXA pipeline as implemented in our codebase
(Section~\ref{sec:mexa-pipeline}), then describe the evaluated models
(Section~\ref{sec:models}), the parallel corpora and their subsets
(Section~\ref{sec:datasets}), and finally the experimental settings,
analysis methods, and compute environment
(Section~\ref{sec:experimental-settings}). The methodology follows
\citet{kargaran2025mexa} for the reproduction experiments and extends it
along four axes: new decoder families, mixture-of-experts architectures,
encoder and embedding models, and non-English pivot languages.

% ------------------------------------------------------------
\section{The MEXA Pipeline}
\label{sec:mexa-pipeline}

MEXA (Multilingual Evaluation via Cross-Lingual Alignment) estimates the
multilingual coverage of a language model $m$ from parallel sentences
alone, in four stages: sentence-embedding extraction, layer-wise
similarity computation, alignment scoring, and aggregation.

\paragraph{Stage 1: Layer-wise sentence embeddings.}
Let $\mathcal{S} = \{s_1, \dots, s_n\}$ be a set of $n$ parallel sentences
available in a target language $L_1$ and a pivot language $L_2$ (English,
unless stated otherwise). Each sentence is passed through the model, and
hidden states $h_t^{(l)}$ are collected for every token $t$ at every layer
$l \in \{1, \dots, L\}$, including the embedding layer. Token states are
reduced to a single sentence embedding $e^{(l)}$ per layer with one of two
strategies:

\begin{itemize}
    \item \textbf{Position-weighted average.} Because causal attention
    means early tokens see little context while late tokens see all of it,
    tokens are averaged with weights proportional to their position:
    \begin{equation}
        e^{(l)} \;=\; \sum_{t=1}^{T} w_t\, h_t^{(l)},
        \qquad
        w_t \;=\; \frac{t}{\sum_{k=1}^{T} k},
        \label{eq:weighted-avg}
    \end{equation}
    where $T$ is the number of non-padding tokens. This is the primary
    setting in all reported results.
    \item \textbf{Last token.} The hidden state of the final non-padding
    token, $e^{(l)} = h_T^{(l)}$, which under causal attention is the only
    position that has attended to the entire sentence. This setting is
    used as a robustness check.
\end{itemize}

For encoder-only models (Section~\ref{subsec:encoder-models}), whose
attention is bidirectional, the same two reductions are applied to the
encoder hidden states, making the pipeline architecture-agnostic.
Embeddings are stored per language and per layer, so that scoring is
decoupled from (and far cheaper than) extraction.

\paragraph{Stage 2: Similarity matrix.}
For each layer $l$, a cosine similarity matrix
$C(L_1, L_2, m, l) \in \mathbb{R}^{n \times n}$ is computed between the
$n$ target-language embeddings and the $n$ pivot-language embeddings:
$c_{ij}(l)$ is the cosine similarity between target sentence $i$ and pivot
sentence $j$. The diagonal entries $c_{ii}(l)$ correspond to true parallel
pairs; all off-diagonal entries are non-parallel distractors.

\paragraph{Stage 3: Alignment score.}
As discussed in Section~\ref{subsec:alignment-challenges}, absolute cosine
values are unreliable in anisotropic LLM embedding spaces \citep{rajaee2022isotropy,hammerl2023anisotropy}.
MEXA therefore scores alignment with a \emph{bidirectional retrieval} (precision-at-1)
criterion: a sentence pair counts as aligned only if its true similarity
strictly exceeds every distractor in both its row and its column,

\begin{equation}
    \mu C(L_1, L_2, m, l)
    \;=\;
    \frac{1}{n}
    \sum_{i=1}^{n}
    \mathbb{1}\!\left(
        c_{ii}(l) \;>\; \max_{j \neq i}\,\{\, c_{ij}(l),\; c_{ji}(l) \,\}
    \right).
    \label{eq:mexa}
\end{equation}

The score lies in $[0, 1]$; the expected value under random embeddings is
approximately $1/n^2$ per pair (both the row and the column condition must
hold by chance), so with $n = 100$ sentences a random baseline scores
essentially zero, which validates the use of small parallel samples.

\paragraph{Stage 4: Aggregation across layers.}
The layer-wise scores are summarized into two model-level statistics:
\begin{equation}
    \mu_{\text{Max}}(L_1) = \max_{l} \; \mu C(L_1, L_2, m, l),
    \qquad
    \mu_{\text{Mean}}(L_1) = \frac{1}{L} \sum_{l=1}^{L} \mu C(L_1, L_2, m, l).
\end{equation}
$\mu_{\text{Max}}$ captures the peak alignment a model achieves anywhere
in its depth --- typically in the middle-to-late layers, consistent with
the pivot hypothesis --- while $\mu_{\text{Mean}}$ rewards models that
sustain alignment across many layers. Model-level scores reported in
Chapter~4 average these quantities over all evaluated languages.

\paragraph{Extension: mean-centered scoring.}
To separate semantic alignment from the language-identity signal
(Section~\ref{subsec:alignment-challenges}), we additionally compute a
\emph{centered} variant of Equation~\eqref{eq:mexa}: before building the
similarity matrix at layer $l$, each language's embeddings are shifted by
their own mean,
\begin{equation}
    \tilde{e}_i^{(l)} \;=\; e_i^{(l)} - \frac{1}{n}\sum_{k=1}^{n} e_k^{(l)},
\end{equation}
applied independently to the target and pivot sides. Centering removes the
dominant direction shared by all sentences of a language and typically
sharpens retrieval; comparing centered and uncentered scores quantifies
how much measured alignment is carried by language-mean offsets rather
than sentence-level semantics.

\paragraph{Extension: non-English pivots.}
Equation~\eqref{eq:mexa} does not privilege English mathematically; the
pivot $L_2$ is a free parameter. We exploit this to test the pivot
hypothesis directly, recomputing all alignment scores with French, German,
Arabic, Chinese, and Basque as pivot languages --- spanning high-resource
Latin-script, non-Latin-script, and a language isolate --- in both the
standard and centered variants.

% ------------------------------------------------------------
\section{Models}
\label{sec:models}

We evaluate 25+ publicly available checkpoints, grouped into four
categories that correspond to the four experimental axes of the thesis.
All models are obtained from the Hugging Face Hub and evaluated without
any fine-tuning; base (non-instruction-tuned) variants are used where
available so that scores reflect pre-training rather than post-training.

% ..............................................
\subsection{Reproduction Models}
\label{subsec:reproduction-models}

To validate our implementation against \citet{kargaran2025mexa}, we
re-evaluate two models with published reference scores:
\textbf{Llama~3.1~8B} \citep{dubey2024llama3}, the primary English-centric
decoder studied in the original paper, and \textbf{Mistral~7B~v0.3}
\citep{jiang2023mistral}. Agreement with the published
$\mu_{\text{Max}}$/$\mu_{\text{Mean}}$ values on the Table-1
configurations (Section~\ref{sec:datasets}) establishes that our
embedding-extraction and scoring code reproduces the original pipeline,
and anchors all subsequent extensions.

% ..............................................
\subsection{New Decoder Families}
\label{subsec:new-decoders}

The first extension applies MEXA to decoder families released after the
original study, covering different data mixtures and scales:

\begin{itemize}
    \item \textbf{Qwen3} \citep{yang2025qwen3} at 0.6B, 1.7B, 4B, and 8B
    parameters --- a family trained with a substantially more multilingual
    data mixture than Llama, whose within-family size sweep isolates the
    effect of scale under a fixed recipe;
    \item \textbf{Qwen3.5~9B}, the successor generation, to measure
    recipe-level progress at approximately constant parameter count;
    \item \textbf{Apertus~8B} and \textbf{Apertus-Mini~4B}, openly trained
    models with an explicitly multilingual, compliance-focused data
    policy, providing a counterpoint to English-centric training;
    \item \textbf{Gemma~4} in five variants (E2B, E4B, 12B, 31B, and the
    sparse 26B-A4B), Google's current open-weight family, spanning
    efficient edge-scale to 31B-parameter models.
\end{itemize}

\noindent
Together with the reproduction models, this yields size ladders within
families (Qwen3 0.6B$\to$8B, Gemma~4 E2B$\to$31B) and recipe contrasts
across families at matched scale (Llama~3.1~8B vs.\ Qwen3~8B vs.\
Apertus~8B).

% ..............................................
\subsection{Mixture-of-Experts Models}
\label{subsec:moe-models}

Sparse mixture-of-experts (MoE) architectures route each token through a
small subset of expert feed-forward networks, decoupling parameter count
from per-token compute \citep{jiang2024mixtral}. Whether expert routing
helps or hurts cross-lingual alignment is an open question: experts could
specialize by language (fragmenting the shared space) or by function
(leaving alignment intact). We evaluate \textbf{Mixtral~8x7B} and
\textbf{Mixtral~8x22B} \citep{jiang2024mixtral}, comparing them against
their dense sibling Mistral~7B~v0.3, and \textbf{Gemma~4~26B-A4B},
comparing it against dense Gemma~4 variants. In all cases the MEXA
pipeline applies unchanged, since it only consumes the post-layer hidden
states.

% ..............................................
\subsection{Encoder and Embedding Models}
\label{subsec:encoder-models}

The final model axis contrasts causal decoders with bidirectional
encoders and purpose-built sentence-embedding models, which provide an
upper-bound reference for what alignment looks like when it is an explicit
training objective rather than a by-product:

\begin{itemize}
    \item \textbf{Masked-LM encoders:} XLM-R base and large
    \citep{conneau2020xlmr}, Glot500 \citep{imani2023glot500} (an XLM-R
    variant extended to 500+ languages, particularly relevant for the
    Bible corpus tail), and mmBERT-base;
    \item \textbf{Contrastively trained sentence encoders:} LaBSE
    \citep{feng2022labse} and multilingual-E5-base
    \citep{wang2024multilingual}, both explicitly optimized so that
    translations are nearest neighbors;
    \item \textbf{Decoder-based embedding models:} Qwen3-Embedding at
    0.6B, 4B, and 8B, which share the Qwen3 backbone but are fine-tuned
    for retrieval --- allowing a controlled comparison of the same
    architecture before and after embedding-specific training.
\end{itemize}

% ------------------------------------------------------------
\section{Datasets}
\label{sec:datasets}

All experiments draw on the two parallel corpora introduced in
Section~\ref{sec:parallel-corpora}, in configurations chosen to match the
original paper where reproduction is intended and to extend coverage
elsewhere.

\paragraph{FLORES-200.}
From the FLORES-200 devtest split \citep{nllbteam2022} we use:
\begin{itemize}
    \item \textbf{FLORES Table-1 subset:} the 116 languages overlapping
    with Belebele, with $n = 100$ parallel sentences per language ---
    the exact configuration of Table~1 in \citet{kargaran2025mexa}, used
    for reproduction and for all model-family comparisons;
    \item \textbf{FLORES Table-1 (2000):} the same 116 languages with all
    $n = 2{,}000$ devtest sentences, to quantify the effect of sample
    size on the retrieval criterion (larger $n$ means more distractors
    and therefore a strictly harder task);
    \item \textbf{FLORES Full:} all 204 language--script pairs, for
    maximal high-quality coverage.
\end{itemize}

\paragraph{Parallel Bible Corpus (sPBC).}
From the super-parallel Bible corpus \citep{mayer2014bible} we use:
\begin{itemize}
    \item \textbf{Bible Table-1 subset:} 101 languages with $n = 103$
    parallel verses, matching the original paper's Bible configuration;
    \item \textbf{Bible Full:} the complete corpus of more than 1{,}400
    language--script pairs, which extends evaluation far into the
    low-resource tail where no downstream benchmark exists.
\end{itemize}

English (\texttt{eng\_Latn}) serves as the pivot in all standard runs; the
pivot-substitution runs of Section~\ref{sec:mexa-pipeline} use the FLORES
Table-1 subset. Language identifiers follow the FLORES-200 convention of
ISO~639-3 code plus script (e.g.\ \texttt{shn\_Mymr}).

% ------------------------------------------------------------
\section{Experimental Settings}
\label{sec:experimental-settings}

% ..............................................
\subsection{Evaluation and Analysis Methods}
\label{subsec:analysis-methods}

Unless stated otherwise, reported scores use position-weighted average
embeddings (Equation~\eqref{eq:weighted-avg}) and are given as
$\mu_{\text{Max}}$ and $\mu_{\text{Mean}}$ per model--dataset
configuration; last-token embeddings and the centered variant are reported
as ablations. The analysis proceeds on four levels:

\begin{itemize}
    \item \textbf{Model level:} scores averaged over all languages of a
    configuration, enabling family- and scale-comparisons and the
    reproduction check against published values;
    \item \textbf{Language level:} per-language scores, ranked and grouped
    by script and by resource level, and cross-referenced with tokenizer
    fertility (computed on both corpora) to separate tokenization effects
    from representational effects;
    \item \textbf{Layer level:} full layer-wise alignment profiles
    $\mu C(\cdot, l)$, visualized as heatmaps, to locate where in the
    network alignment peaks and to test the intermediate-layer prediction
    of the pivot hypothesis;
    \item \textbf{Validity level:} Pearson correlations between MEXA
    scores and downstream accuracy on Belebele
    \citep{bandarkar2024belebele} and m-ARC \citep{lai2023okapi} over the
    languages both cover, testing whether alignment is a faithful proxy
    for task performance; cross-model Pearson correlations of per-language
    score profiles additionally measure how consistent alignment
    structure is across architectures.
\end{itemize}

For qualitative inspection, sentence embeddings are also projected to two
and three dimensions with PCA and t-SNE, colored by language and script,
to visualize the clustering behavior that the retrieval score summarizes
numerically.

% ..............................................
\subsection{Implementation and Compute}
\label{subsec:implementation}

The pipeline is implemented in Python with PyTorch and Hugging Face
Transformers. \texttt{embed\_extractor.py} loads each checkpoint via
\texttt{AutoModelForCausalLM} (or \texttt{AutoModel} for encoders) with
output of all hidden states enabled, tokenizes each sentence with the
model's own tokenizer, and stores per-layer sentence embeddings as
serialized NumPy arrays, one file per language.
\texttt{compute\_mexa.py} then builds the layer-wise cosine similarity
matrices in float64 and evaluates Equation~\eqref{eq:mexa}; because
scoring operates on cached embeddings, ablations over pooling strategies,
pivots, and centering re-use a single (expensive) extraction pass per
model--corpus pair. Formatted per-language and per-layer scores are
exported to CSV/JSON and explored in an interactive React dashboard
developed for this thesis.

Small models were run locally; all large-scale extraction jobs (8B and
above, MoE, and full-corpus sweeps) were executed on the CIT-TUM-HN
SLURM cluster at TUM Campus Heilbronn on NVIDIA data-center GPUs, with
models loaded in reduced precision (bfloat16, with quantized loading for
the largest MoE checkpoints) and embeddings cast to float32 before
storage. Random effects play no role in the pipeline --- extraction and
scoring are deterministic given a checkpoint and corpus --- so no seed
averaging is required.
